% Methods section -- coupled-oscillator causal discovery with order-parameter conditioning.
% Requires: \usepackage{amsmath,amssymb}  (and a bibliography for the \cite commands).
% Drop into the manuscript after the Introduction.

\section{Methods}
\label{sec:methods}

\subsection{Coupled phase-oscillator model and emergent regimes}
\label{sec:model}

We consider $N$ coupled phase oscillators evolving under a generalized
Kuramoto--Sakaguchi system,
\begin{equation}
\dot\theta_i \;=\; \omega_i \;+\; K\sum_{j=1}^{N}\widetilde{W}_{ij}\,
\Gamma\!\left(\theta_j-\theta_i\right)\;+\;\xi_i(t),
\qquad
\Gamma(\phi)=\sin(\phi-\alpha)+a_2\sin\!\big(2\phi-\beta\big),
\label{eq:model}
\end{equation}
where $\theta_i$ is the phase of oscillator $i$, $\omega_i$ its natural
frequency, $K$ the coupling strength, and $\xi_i$ independent Gaussian phase
noise of standard deviation $\sigma$. The coupling function $\Gamma$ carries a
phase lag $\alpha$ and an optional second harmonic of amplitude $a_2$ and lag
$\beta$; $\widetilde{W}_{ij}=W_{ij}/\sum_k W_{ik}$ is the row-normalized coupling
matrix. Equation~\eqref{eq:model} is integrated with a fourth-order
Runge--Kutta scheme. A directed coupling $j\!\to\! i$ exists iff $W_{ij}>0$, and
the binary support $A=\operatorname{supp}(W)$ is the \emph{ground-truth causal
graph} to be reconstructed. By selecting the topology $W$, the lag $\alpha$, the
harmonic content $(a_2,\beta)$, and the frequency distribution, Eq.~\eqref{eq:model}
is driven into four emergent regimes---chimera, traveling-wave, metastable, and
cyclops states (Table~\ref{tab:regimes})---each verified \emph{a posteriori},
after discarding a transient, by the order-parameter diagnostics of
Sec.~\ref{sec:orderparam}.

\subsection{Collective coordinates: the order parameters}
\label{sec:orderparam}

The collective state of the ensemble is summarized by the complex
Kuramoto--Daido order parameters
\begin{equation}
Z_m(t)\;=\;R_m(t)\,e^{i\psi_m(t)}\;=\;\frac{1}{N}\sum_{j=1}^{N}e^{\,i m\theta_j(t)},
\qquad m=1,2,
\label{eq:Zm}
\end{equation}
with $R_m\in[0,1]$ the coherence and $\psi_m$ the collective phase of the
$m$-th harmonic. We use $Z_1,Z_2$ both to certify emergence (e.g.\ a
traveling wave by $|Z_1|\!\to\!0$ with $|Z_q|\!\approx\!1$ for the twisted
mode $q$; a cyclops state by its cluster partition and the second-harmonic
coherence $R_2$) and, crucially, as conditioning variables for causal
discovery (Sec.~\ref{sec:opcond}).

\subsection{Physics-informed basis and optimal causation entropy}
\label{sec:ocse}

Causal reconstruction is performed on a physics-informed linearization of
Eq.~\eqref{eq:model}. For each oscillator $i$ we take the one-step phase
velocity $v_i(t)=\big[\theta_i(t{+}dt)-\theta_i(t)\big]/dt$ as the regression
target and, as candidate predictors, the pairwise coupling features
\begin{equation}
s_{i\leftarrow j}(t)\;=\;\sin\!\big(\theta_j(t)-\theta_i(t)\big),
\qquad j\neq i,
\label{eq:basis}
\end{equation}
which are exactly the terms entering the equation of motion. Selecting the
feature $s_{i\leftarrow j}$ for target $i$ corresponds to the directed edge
$j\!\to\! i$.

Optimal causation entropy (oCSE) \cite{Sun2015} identifies the parents of each
target by aggregative discovery and removal on the \emph{causation entropy}
\begin{equation}
C_{\,j\to i\,\mid\,\mathcal{Z}}
\;=\; I\!\left(v_i^{\,t+1}\,;\,s_{i\leftarrow j}^{\,t}\;\middle|\;\mathcal{Z}^{\,t}\right),
\label{eq:ce}
\end{equation}
the conditional mutual information between the target's future and a candidate
feature given a conditioning set $\mathcal{Z}$. A feature is admitted only if
$C_{\,j\to i\,\mid\,\mathcal{Z}}$ is significantly positive and is later removed
if it becomes redundant given the selected set, so that $\mathcal{Z}$ grows to
contain the target's own past and its already-selected parents. The baseline
estimator conditions on the target's own history,
$\mathcal{Z}=\{v_i^{\,t}\}\cup\{\text{selected parents}\}$.

\subsection{Conditioning on the order parameter}
\label{sec:opcond}

The central methodological element of this work is to augment the oCSE
conditioning set with the collective coordinates of
Eq.~\eqref{eq:Zm}. We define the \emph{order-parameter conditioning vector}
\begin{equation}
O(t)=\big(\operatorname{Re}Z_1,\ \operatorname{Im}Z_1,\
\operatorname{Re}Z_2,\ \operatorname{Im}Z_2\big)(t),
\label{eq:Ovec}
\end{equation}
and run oCSE with $O(t)$ permanently in the conditioning set,
$\mathcal{Z}=\{v_i^{\,t}\}\cup O(t)\cup\{\text{selected parents}\}$; $O$ is
supplied as a conditioning variable only and is never itself selectable as an
edge. We refer to this estimator as oCSE+OP. The rationale is that, in
synchronized and partially synchronized regimes, the order parameter is
precisely the variable through which the ensemble exerts a \emph{common} drive
on every oscillator, and is therefore the confounder responsible for
synchrony-induced false positives.

This is made precise by the mean-field reduction of Eq.~\eqref{eq:model}. Writing
the coupling through the order parameters and using
$\tfrac1N\sum_j e^{im\theta_j}=Z_m$, the population term acting on oscillator $i$
becomes
\begin{equation}
\frac{1}{N}\sum_{j=1}^{N}\Gamma(\theta_j-\theta_i)
=\operatorname{Im}\!\big[e^{-i(\theta_i+\alpha)}Z_1\big]
+a_2\,\operatorname{Im}\!\big[e^{-i(2\theta_i+\beta)}Z_2\big],
\label{eq:meanfield}
\end{equation}
so that, in the mean-field (globally coupled) limit, the equation of motion
collapses to
\begin{equation}
\dot\theta_i=\omega_i+K\Big[R_1\sin(\psi_1-\theta_i-\alpha)
+a_2R_2\sin(\psi_2-2\theta_i-\beta)\Big]
\;\equiv\; g\!\big(\theta_i,\omega_i,Z_1,Z_2\big).
\label{eq:greduction}
\end{equation}
Equation~\eqref{eq:greduction} states that each oscillator's increment depends on
the \emph{entire rest of the ensemble only through} $(Z_1,Z_2)$. The pair
$(Z_1,Z_2)$ is thus a \emph{sufficient statistic} for the collective influence:
it is the complete common cause shared by all oscillators. Consequently, for any
non-adjacent pair, the candidate feature carries no information about the
target's future once the order parameter is given,
\begin{equation}
I\!\left(v_i^{\,t+1}\,;\,s_{i\leftarrow k}^{\,t}\;\middle|\;\theta_i^{\,t},\,Z_1^{\,t},\,Z_2^{\,t}\right)=0
\qquad\text{whenever } A_{ik}=0,
\label{eq:dsep}
\end{equation}
because oscillator $k$ enters $v_i^{\,t+1}$ exclusively through the collective
coordinates. In graphical terms, $O(t)$ blocks the back-door path
$s_{i\leftarrow k}\!\leftarrow\! Z \!\rightarrow\! v_i$ through which global
synchrony manufactures dependence between unconnected oscillators; conditioning
on $O$ $d$-separates them. A marginal or pairwise estimator (e.g.\ Granger
causality) measures exactly this synchrony-induced dependence and reports it as a
spurious edge, whereas the causation entropy of Eq.~\eqref{eq:ce} with $O$ in the
conditioning set vanishes for non-edges and they are correctly rejected.

For networks with sparse $W$ the reduction is exact only for the \emph{local}
field over a node's neighborhood; the global order parameter $O(t)$ then captures
the dominant \emph{common mode}, which is precisely the component shared across
all nodes and hence the source of confounding in (partially) synchronized
states. What survives conditioning on $O$ is the residual, oscillator-specific
coupling---non-zero only for genuine neighbors, whose feature $s_{i\leftarrow j}$
enters $v_i$ as an explicit additive term beyond the mean field. Conditioning on
the order parameter therefore removes the synchrony-induced false positives while
preserving the network-specific structure. The same conditioning can, however,
suppress a weak true edge whose contribution is nearly collinear with the
collective drive; oCSE+OP is thus a precision-favoring estimator that trades a
small amount of sensitivity for a substantial reduction in false positives, a
trade-off we quantify in Sec.~\ref{sec:results}. The second harmonic $Z_2$ is
included because biharmonic coupling [the $a_2$ term in Eq.~\eqref{eq:model}, as
in cyclops states] makes the collective drive two-mode, so that
Eq.~\eqref{eq:greduction} requires both $Z_1$ and $Z_2$ for sufficiency.

\subsection{Estimation, baselines, and evaluation}
\label{sec:estimation}

Causation entropies in Eq.~\eqref{eq:ce} are estimated under a Gaussian
assumption (partial correlation), and significance is assessed by a permutation
(shuffle) test at level $\alpha$ with $S$ surrogates; we use $S\ge20$ so that the
attainable $p$-value resolution $1/(S{+}1)$ lies below $\alpha=0.05$. We compare
oCSE and oCSE+OP against PCMCI (with a partial-correlation conditional
independence test) \cite{Runge2019}, pairwise linear Granger causality, and
VARLiNGAM \cite{Hyvarinen2010}, each applied to the same basis of
Eq.~\eqref{eq:basis} and reduced to a node-level adjacency by the same
feature-to-edge map. Reconstructed graphs are scored against
$A=\operatorname{supp}(W)$ by true- and false-positive rate, precision, and
$F_1$. Because the cyclops coupling is all-to-all, its ground-truth graph is
complete and admits no true negatives; there the false-positive rate is undefined
and only recall and precision are reported.
